% =====================================================================================================
\section{Discussão}\label{sec:discussao}
\todo{síntese por família; efeito das iscas e do desperdício (dano $\times$ custo); efeito da composição e do voo
baixo sobre a previsão de velocidade constante; generalização para cidades nunca vistas; custo computacional e
escalabilidade; recomendações.}

\subsection{Limitações}
\begin{itemize}
  \item \textbf{Modelo de interceptador.} Os interceptadores são pontos de massa em duas dimensões, com velocidade
        constante e aceleração lateral limitada; não há dinâmica de atitude nem vento sobre eles. Por isso, a altura da
        ameaça afeta a detecção pelo radar, mas não o engajamento.
  \item \textbf{Modelo de ameaça.} A velocidade de cada ameaça é constante ao longo da rota (o aumento de velocidade da
        munição vagante no mergulho não é modelado), e só o trecho final das rotas de longo alcance, dentro de um mapa de
        12 a 14~km, é simulado. A seção radar, o valor de dano de cada classe e a penalidade por desperdício
        ($\eta=0{,}2$) são hipóteses do modelo; os parâmetros de fontes abertas vêm de reportagens e análises
        secundárias, não de especificações oficiais.
  \item \textbf{Ativos.} Os alvos dependem do mapeamento colaborativo do OpenStreetMap, cuja cobertura varia entre as
        cidades; a importância de cada categoria de ativo é uma escolha do modelo.
  \item \textbf{Sensor próprio ideal.} Na fase final, o interceptador vê a aeronave sem erro.
  \item \textbf{Adaptações.} As técnicas foram adaptadas a interceptadores descartáveis de contato único. A árvore com
        utilidade (M23) perde o ramo de reforço, e o MILP (M15) representa só a parte de alocação de~\cite{cruz2025}, sem
        obstáculos nem bateria. Só as redes treinadas por reforço conhecem a penalidade por desperdício; os métodos de
        WTA usam o seu próprio custo de lançamento, e os demais não têm nenhum.
  \item \textbf{Modelo de sucesso.} Os métodos M13--M20 dependem de~\eqref{eq:q}, calibrado num experimento 1$\times$1; um
        modelo melhor poderia favorecer essa família.
  \item \textbf{Orçamento de treino.} As redes foram treinadas com um orçamento fixo e modesto; resultados com mais
        treino podem diferir.
\end{itemize}

\section{Conclusão}\label{sec:conclusao}
\todo{}

\section*{Disponibilidade de código e dados}
O simulador, os geradores de missões, a extração dos ativos do OpenStreetMap, as 192 missões, as implementações dos 29
métodos (um arquivo por método), os modelos treinados e os scripts que geram as tabelas e figuras deste artigo estão
disponíveis em \todo{repositório}. Qualquer método pode ser rodado isoladamente sobre qualquer missão
(\texttt{rodar.py}).
